\documentclass[10pt]{article}
\usepackage[margin=0.8in]{geometry}
\usepackage[]{xcolor}
\usepackage{amsmath}
\parindent=0pt

\title{Exam 2 Review}
\author{gordonbchen}
\date{}

\begin{document}
\maketitle


\section{Ch 3}
Starting at ``design issues for paging system"

\begin{enumerate}
    \item Local vs global alloc

        local: page replaces page in same process (working set is local)

        global: page replaces page in any process

    \item space per proc: equal alloc, proc size

        page fault rate per proc, boundary for page fault rate, high page fault rate = get more pages

        load control: all procs need more pages, suspend some and give up pages, mix CPU-bound + IO-bound

    \item page table overhead = (process size / page size) * page entry size + (page size / 2)

        optimal page size = sqrt(2 * process size + page entry size)

    \item separate instruction and data spaces

        same program running multiple times: instruction space is the same, data space is copied on write

        shared libraries (dll + so): programs have pointer to library in ram (don't have to load multiple times, update at once)

    \item cleaning policy

        paging daemon: periodically evict + write pages if available pages is too low

        2 handed clock: 1st hand = cleaning daemon, 2nd hand = page replacement alg

    \item page fault: hardware interrupt (not in TLB), kernel checks page table + maybe fetch page

    \item instruction backup: pipe

    \item lock pages in memory: A does IO + gets suspended, B gets page fault (global page replacemnet), claims file A is using for IO

        pinning: lock page

        kernel buffers for IO

    \item swap: store pages not in page table in disk TODO

    \item segmentation: each variable can get own segment (data, stack, heap), programmer defines segments

        without paging: whole segments get brought in / swapped out, compaction (address translation w/ + base register), physical addr = segment table[segment number] + offset

        with paging: MULTICS
        \begin{enumerate}
            \item descriptor segment has pointers to page tables
            \item multics virtual address = segment number + page number + offset
            \item physical addr = segment table [segment number] [page number] + offset
            \item multics tlb: index on segment number, virtual addr, gets page number
        \end{enumerate}

        Pentum: 256K segments, 64k words
        \begin{enumerate}
            \item global (0) (for OS / system) vs local (1) descriptor table (for program)
            \item selector: which segment, offset: which word in segment
            \item selector gets global/local descriptor, extract segment base addr, linear addr = segment base + offset

                without paging: linear addr = physical addr

                with paging: linear addr = virtual addr = page dir number | page table number | offset
        \end{enumerate}

        Intel x86: 16k segments, 1 billion words

        x86-64: no segmentation
\end{enumerate}


\section{Ch 6: Deadlocks}
\begin{enumerate}
    \item detection + recovery: graph / matrix, TODO
    \item dynamic avoidance: check at resource allocation time
    \item prevention:

        mutual exclusion: not practical b/c somethings need single access

        circular wait: number resources in increasing order, only can request resources adter you

    \item 2 phase locking: databases

        communication deadlock: timeout + sequence numbering

        livelock: dining philosophers

        priority starvation
\end{enumerate}


\section{Ch. 4: Filesystems}
\begin{enumerate}
    \item file system layout
        \begin{enumerate}
            \item bios, mbr, boot block, loads os
            \item master boot record: which is active partition
            \item partitions: boot block ...
        \end{enumerate}

    \item allocation

        contiguous

        linked list alloc: pointer to next, block. no random access, must be sequential

        FAT (file allocation table): proportional to size of block

        INodes (linux): file attr + addr to data blocks + addr to (1x, 2x, 3x) indirect blocks
        \begin{enumerate}
            \item indirect block: pointer to block of pointers to blocks
        \end{enumerate}

    \item backups

        physical dump: dump everything

        logical dump: bitmap for modified files, only dump chains of dirs w/ modified files
\end{enumerate}

\end{document}
